\section{Fixing the Terms}

The three words belong to different orders. One names a virtue, one names a behaviour, one names a status in law. Nothing in what follows works unless that is held steadily, so it is worth establishing each from the texts that define it rather than from usage.

\subsection{Chastity: the virtue that binds everyone}

Chastity is not a clerical obligation at all, in the sense that it is not an obligation \emph{of} clerics as distinct from anyone else. The \emph{Catechism} states the rule and its universality in one line: \latin{Omnis baptizatus ad castitatem vocatur}---all the baptized are called to chastity---and immediately adds the qualifier that does the real work: ``All Christ's faithful are called to lead a chaste life in keeping with their particular states of life.''\endnote{CCC 2348, English quoted from the Holy See's archival English delivery of the sixth-commandment article, fetched, hashed, and read 2026-07-25; Latin from the Holy See's Latin delivery of the same article (\texttt{catechism\_lt/p3s2c2a6\_lt.htm}), fetched, hashed, and read the same day. The Latin typical edition governs; the English is an identified official delivery used as a working aid. See References for both locations.} The obligation is universal in its subject and differentiated in its content. The next paragraph spells the differentiation out, quoting the 1975 declaration \emph{Persona humana}: some ``profess virginity or consecrated celibacy''; others ``live in the way prescribed for all by the moral law, whether they are married or single''; and then the sentence that separates our second and third terms from the first---``Married people are called to live conjugal chastity; others practice chastity in continence,'' \latin{ceterae castitatem colunt in continentia}.\endnote{CCC 2349, same deliveries; the internal quotation is from the Congregation for the Doctrine of the Faith, declaration \emph{Persona humana} (29 December 1975), n.~11, which the \emph{Catechism} cites at that place. \emph{Persona humana} was not separately examined for this article; the quotation is used as the \emph{Catechism} carries it, which is the level at which the claim here is made.}

So chastity is a genus with species, and the species are assigned by state of life. A married man practises chastity in and through conjugal life; he is chaste precisely by being a faithful husband, not despite it. An unmarried man practises chastity by continence. This is why the Eastern Code can bind celibate and married clerics alike to \latin{castitatis decor} in a single clause without saying anything at all about who may marry: chastity is the common term. It is also why ``a chaste priesthood'' is not a description of the Latin discipline but of the whole Church, and why arguing for clerical celibacy from the value of chastity proves nothing, since chastity is not what is at issue between the disciplines.

Two further distinctions belong here. Chastity is a \emph{virtue}, and canon law does not legislate virtues directly; it legislates acts, states, and conditions. When a code speaks of chastity it is either exhorting (CCEO c.~374's \latin{elucere debent}), or naming a counsel that generates a determinate juridical obligation (CIC c.~599), or penalizing an external act (CIC c.~1395 §1, which punishes the cleric who persists ``in another external sin against the sixth commandment of the Decalogue with scandal'').\endnote{CIC c.~1395 §1 in the text of Book~VI as integrally revised by \emph{Pascite gregem Dei} (23 May 2021, in force 8 December 2021), Latin delivery of Book~VI fetched, hashed, and read 2026-07-25 and byte-identical to the state registered in this repository's source library: \latin{Clericus concubinarius, praeter casum de quo in can.~1394, et clericus in alio peccato externo contra sextum Decalogi praeceptum cum scandalo permanens, suspensione puniantur\ldots} Note the canon's own limits: it reaches the \emph{external} and \emph{scandalous} and persistent, which is what a penal law can reach; the virtue itself is not its object.} The virtue is presupposed everywhere and adjudicated nowhere.

\subsection{Continence: the abstention}

Continence is a behaviour: abstention from sexual activity. It is what canon 277 §1 obliges, in a formula the Code repeats for religious at canon 599 and the Second Vatican Council used at \emph{Presbyterorum ordinis} 16---\latin{perfecta et perpetua propter Regnum coelorum continentia}, perfect and perpetual continence for the sake of the Kingdom of heaven.\endnote{\emph{Presbyterorum ordinis} 16, Latin quoted from the Holy See's Latin web delivery of the Decree, fetched, hashed, and read 2026-07-25 and byte-identical to the state registered in this repository's source library; the article's Latin and English deliveries were previously registered from 2026-07-23 fetches. The formula's three components are worth separating: \emph{perfecta} (not partial), \emph{perpetua} (not for a season), \latin{propter Regnum coelorum} (the motive, drawn from Matthew 19:12, which the Council and the Code both retain in the operative sentence rather than relegating to a preamble).} The two adjectives are the whole content: \emph{perfect}, so that no sexual activity is excepted, and \emph{perpetual}, so that no period is excepted. A temporary or partial continence---the continence of Lent, or of spouses practising periodic abstinence---is continence in the same sense and neither perfect nor perpetual.

Continence and chastity are not the same and can come apart in both directions. An unmarried man who abstains entirely from sexual acts but lives in habitual lust is continent and unchaste; a married man who has never abstained a day in his life may be entirely chaste. The law tracks this. What canon 277 §1 obliges is the abstention, which is externally determinate and therefore a fit object of legal obligation; the virtue that should animate it is the object of the exhortation in §2 and of formation, not of a norm with a sanction.

\subsection{Celibacy: the state}

Celibacy is neither a virtue nor a behaviour but a status: the condition of being unmarried. In canon law it is a status with three distinct juridical faces, and it is worth naming them separately because they are established by three different canons and can be undone by three different acts. It is (i) an \emph{obligation} publicly assumed before ordination (c.~1037); (ii) a \emph{diriment impediment} to marriage attaching to those in sacred orders (c.~1087); and (iii) a condition whose breach is a \emph{delict} carrying a \latin{latae sententiae} penalty (c.~1394 §1). The obligation can be dispensed only by the Roman Pontiff (c.~291); the impediment can be dispensed only by the Apostolic See (c.~1078 §2 1°); the penalty is remitted by the competent authority under the ordinary penal norms. Sections~4 through~6 work each of these through.

Note what celibacy is not. It is not continence, though in the Latin discipline it is imposed \emph{because} of continence---the \latin{ideoque} of canon 277 §1 runs from the continence to the celibacy, not the reverse. A widowed permanent deacon who has never remarried is celibate as a matter of status; whether he is thereby obliged to perfect continence is a question the Latin Code raises and does not clearly answer, which is the subject of section~9. And celibacy is not consecration: the celibacy of a diocesan presbyter is not a religious vow, is not a public vow at all, and is governed by the law of clerical obligations rather than by the law of consecrated life. That distinction is why canon 1087 and canon 1088 are separate canons with separate objects, and why canon 1078 §2 1° names the impediment ``arising from sacred orders \emph{or} from a public perpetual vow of chastity in a religious institute of pontifical right'' as two things.

\subsection{The three terms in one view}

\begin{fourstable}{Term}{What it names}{Who is bound, and by what}{How it can cease}
Chastity & The virtue integrating sexuality into the person, differentiated by state of life. & Every baptized person, by the moral law; CCC 2348--2349. Named in law at CIC c.~599 (as counsel) and CCEO c.~374 (as exhortation). & It cannot cease. No state of life, dispensation, or authority exempts anyone from it.\\
Continence & Abstention from sexual activity; ``perfect'' when total, ``perpetual'' when lifelong. & Latin clerics, by CIC c.~277 §1; religious, by CIC c.~599; not imposed as such on Eastern clerics by CCEO common law. & With the obligation that carries it: for a Latin cleric, by papal dispensation from c.~277 §1 (c.~291), which in practice accompanies loss of the clerical state.\\
Celibacy & The unmarried state, considered juridically: obligation, marriage impediment, and penal condition. & Latin clerics: obligation assumed at c.~1037, impediment at c.~1087, penalty at c.~1394 §1. Eastern clerics: impediment at CCEO c.~804; obligation only where particular or special law imposes it. & The obligation, by papal dispensation alone (CIC c.~291; CCEO c.~396). The impediment, by dispensation reserved to the Apostolic See (CIC c.~1078 §2 1°; CCEO c.~795 §1 1°).\\
\end{fourstable}

The right-hand column is where the three terms stop being a matter of vocabulary. One of them admits of no exception whatever; one is dispensable by a single man on earth; one is an impediment whose relaxation is a reserved administrative act. A discipline described in terms that cannot distinguish these cannot be discussed accurately, defended accurately, or criticized accurately.
